\section{Protocol for the preliminary ten-task comparison}\label{app:matched}
Each request provided one task, its central examples and three candidate architectures with their initialized coefficients. The response had to supply one complete array of replacement weights for each architecture. We crossed ten tasks, three random-number seeds (164, 317, 583) and two information conditions, yielding 60 independent calls and 180 arrays. Three arrays in one response share a prompt. A seed changes the initialization and the sampled examples, except that the two canonical number sequences always use the same indices. The \emph{Formula and examples} condition supplied the task definition and central input--output pairs. The \emph{Examples only} condition omitted the formula and task name but retained informative choices of architecture, indexing and deployment range. Every saved command-line interface trace recorded one response and zero tool actions. Each measurement uses the first returned array.

For the four smooth targets, the choices were one, four or 16 scalar hinges with an affine skip from input to output (5, 14 or 50 adjustable scalars). For addition and subtraction, the three models saw only the first input, both inputs without bias, or both inputs with a bias (1, 2 or 3 scalars). For running sum, the models were a memoryless linear map or recurrent networks with one or two hidden coordinates (1, 3 or 8 scalars). For sorting, let \(h\in\mathbb R^4\) denote the current four values and let \(v,w\in\mathbb R^4\) be learned columns. A block updates \(h\) to \(h+v\relu(w^\top h)\), using eight unconstrained coefficients. Four, five or six such blocks use 32, 40 or 48 coefficients and permit different routing in every block. For a sequence generator of width \(d\), the adjustable entries are a \(d\)-by-\(d\) transition, \(d\) initial values and \(d\) readout weights. The tested widths were \(1,2,3\) for Fibonacci and \(2,3,4\) for triangular numbers. Table~\ref{tab:paired} selects the 16-hinge smooth network, the two-input affine map, and the smallest known sufficient real-arithmetic configurations for the remaining tasks.

The WSP and backpropagation conditions received 12 central labeled examples per task, except for the canonical sequence generators, which used the nine indices \(0,\ldots,8\). Input lists had length four for sorting and six for sums. Fixed-length tests held the lengths constant and changed the input magnitudes. A scalar shell side contributed 2,048 inputs, and positive self-power used only the right side. Each vector test comprised 2,048 independent examples. Sums also used separate length-32 tests on central and outer values. Canonical number sequences have an early transfer set at indices 9--32 and a later outer set at 33--64. An exact percentage there counts correct individual numbers along one trajectory. For a vector or stream, exactness requires equality at every coordinate or tick within that example.

Numerical labels come from the stated NumPy binary64 operation on those inputs. The separate smooth-function studies use the GNU MPFR library for correctly rounded multiple-precision arithmetic, which supplies an independent reference for transcendental functions. Training used 800 full-batch updates of adaptive moment estimation (Adam) and up to 1,000 Broyden--Fletcher--Goldfarb--Shanno (BFGS) iterations using an approximate inverse curvature matrix. For each scalar reference \(y\), the residual is the prediction minus \(y\), divided by \(\max(1,|y|)\). Training minimizes half the mean of the squared residuals. Each gradient at initialization was checked by finite differences, comparing it to a symmetric small perturbation of each parameter. The WSP prompt prioritized the largest absolute scaled residual. To compare both methods under one criterion, Appendix~\ref{app:optimizer} scores all returned weights under the same training loss and gives outer root mean squared error (RMSE) in the original output units.
